\documentclass[11pt]{article}
\usepackage[T1]{fontenc}
\usepackage{lmodern,amsmath,amssymb,amsthm,mathtools}
\usepackage[margin=28mm]{geometry}
\usepackage[hidelinks]{hyperref}
\usepackage{microtype}
\setlength{\emergencystretch}{3em}
\newtheorem{theorem}{Theorem}[section]
\newtheorem{lemma}[theorem]{Lemma}
\newtheorem{proposition}[theorem]{Proposition}
\newtheorem{remark}[theorem]{Remark}
\newcommand{\Var}{\operatorname{Var}}
\newcommand{\tr}{\operatorname{tr}}
\newcommand{\E}{\mathcal E}
\newcommand{\Pcal}{\mathcal P}
\title{Rectangular mean observations with critical inverse moments}
\author{Lluis Eriksson}
\date{7 November 2026 --- grouped supplement for the existing paper dossier}
\begin{document}
\maketitle
\begin{abstract}
We construct rectangular mean observations for every positive scalar
coefficient on a finite periodic lattice that depends on one coordinate.
At the critical inverse-moment exponent, the number of cells attains the
target upper order, with the penalty coefficient equal to the coercivity
threshold. The proof uses actual interval resistances, transverse averaging
and a balanced-tree charge valid for odd lattice sizes. No bounded contrast
or separation of weak sites is assumed. A thin weak layer separates this
construction from certificates based only on the minimum coefficient.
We also determine the exact asymptotic cell-cost phase diagram of strong
Green-moment certificates on homogeneous coefficients, for arbitrary
rectangular partitions and aspect ratios. At the critical exponent these
certificates pay a logarithmic overhead although one actual mean suffices.
An explicit coexact field refutes extension of the laminate argument by
sampling a representative fibre. The unrestricted critical rectangular
question remains open. These finite analytical results consolidate related
positive and negative mechanisms; they do not establish a physical mass gap,
formal verification or independent originality.
\end{abstract}

\section{Problem and relation to the existing record}
This grouped supplement belongs to the existing coarse-information paper
dossier~\cite{programme}. It combines the complementary research units184
and185 rather than proposing a deposit for each corollary. Let $D\ge3$,
$n\ge2$, $X=(\mathbb Z/n\mathbb Z)^D$, $N=n^D$ and $p=D/2$.
All differences have unit spacing and all norms use counting measure. Set
\[
 \E_a(u)=\sum_{x\in X}\sum_{j=1}^D
 a(x)|u(x+e_j)-u(x)|^2,\qquad a(x)>0.
\]
For a disjoint rectangular partition $\Pcal$ define the actual normalized
arithmetic means $C_{\Pcal}u=(\sqrt{|R|}\,\bar u_R)_{R\in\Pcal}$.
Our positive construction gives, simultaneously for all fields,
\begin{equation}\label{eq:goal}
 c\|u\|^2\le\E_a(u)+c\|C_{\Pcal}u\|^2,\qquad c>0.
\end{equation}
For a fixed $h$-component field the output dimension is $h\#\Pcal$ over
the selected scalar field; complex dimension must be doubled when counted
over the reals. Scalar estimates sum over components, and hence restrict
to any actual coexact subspace without assuming an abstract surrogate.

The original question asks for a target-order partition under
$\mu=\operatorname{mean}(a)\ge1$, $1\le a\le2\mu$ and
$\operatorname{mean}((\mu/a)^p)\le M$, with cell count bounded by a constant
depending on $D,c,M$ times $\max(1,N/(1+\mu)^p)$. We solve here only the
subclass $a(x)=\alpha(x_1)$. Lamination is an explicit retained hypothesis.
The Green-moment results below concern the cost of a sufficient certificate,
not a lower bound for every successful observation or for coexact rank.

\section{A constructive theorem for arbitrary laminates}
Put $d=D-1$ and
\[
 r=(2/3)^d,\qquad B_D=\frac{2\,6^d}{1-r}.
\]
\begin{theorem}\label{thm:laminate}
For every positive sequence $\alpha(0),\ldots,\alpha(n-1)$, set
$a(x)=\alpha(x_1)$. There is a disjoint axis-aligned rectangular partition
independent of $u$ that satisfies~\eqref{eq:goal} and
\begin{equation}\label{eq:count}
 \#\Pcal\le2^d+B_D(c/2)^p\sum_{x\in X}a(x)^{-p}.
\end{equation}
If the original critical inverse-moment assumptions hold, its output rank
is at most
\[
 h\bigl[2^d+B_D(c/2)^p2^p M\bigr]
 \max\bigl(1,N/(1+\mu)^p\bigr).
\]
The upper bound $a\le2\mu$ is not needed for~\eqref{eq:count}.
Target order here means the requested upper order, not a matching lower
bound for every individual coefficient.
\end{theorem}

\subsection{Actual local energy estimate}
For an interval $I$ of length $m$, write
$R_I=\sum_{x\in I}\alpha(x)^{-1}$. Include its last starting vertex in
this sum: although its longitudinal edge may leave the cut interval, its
coefficient is needed for transverse averaging. Internal rectangle energies
keep only edges within cut paths; periodic closing and cross-cell edges
are discarded as nonnegative contributions.

\begin{lemma}\label{lem:local}
If $R=I\times Z$ and every transverse side has length at most $m$, then
\begin{equation}\label{eq:local}
 \Var_R(u):=\sum_R|u-\bar u_R|^2\le\frac{mR_I}{2}\E_{a,R}(u).
\end{equation}
The constant $mR_I/2$ is an inverse-gap bound.
\end{lemma}
\begin{proof}
Let $v(y)=m^{-1}\sum_{x\in I}u(x,y)$. Orthogonal decomposition gives
\[
 \Var_R(u)=\sum_{y\in Z}\Var_I(u(\cdot,y))+m\Var_Z(v).
\]
The pairwise variance identity and weighted Cauchy along a path give
$\Var_I\le(m-1)R_I\E_1/2\le mR_I\E_1/2$.
For an unweighted path of length $\ell$, the same identity bounds the
coefficient of its $k$th edge by $k(\ell-k)\le\ell^2/4$:
use the path Cauchy bound with its length and then bound that length
divided by $\ell$ by one. Tensorizing these path inequalities gives
$\Var_Z(v)\le m^2\sum_{y,j>1}|D_jv(y)|^2/4$; the tensorized constant
is the largest coordinate constant, not their sum. Weighted Cauchy yields
\[
 |D_jv(y)|^2\le\frac{R_I}{m^2}
      \sum_{x\in I}\alpha(x)|D_ju(x,y)|^2.
\]
Thus $\Var_R\le(mR_I/2)\E_1+(mR_I/4)\E_\perp$, proving the lemma.
This is the step where a common coefficient sequence on all fibres is used.
\end{proof}

\subsection{Partition algorithm and critical charge}
Start from the interval $\{0,\ldots,n-1\}$. Split a node of length $m\ge2$
into consecutive lengths $\lfloor m/2\rfloor,\lceil m/2\rceil$ exactly
when $mR_I>2/c$. Accept equality. Never split a singleton. For each final
interval of length $m$, partition every transverse coordinate into intervals
of length $m$ with a possible shorter last interval. Their products form
an actual disjoint cover. A nonsingleton leaf satisfies
$c\Var_R\le\E_{a,R}$ by Lemma~\ref{lem:local}; singleton longitudinal
leaves produce only singleton cells, with zero variance. Summing and using
$\|u\|^2=\sum_R\Var_R(u)+\sum_R|R||\bar u_R|^2$ proves~\eqref{eq:goal}.

For the count set $b(I)=\sum_I\alpha^{-p}$ and $w(I)=m^{-d}$.
H\"older gives $R_I\le m^{1-1/p}b(I)^{1/p}$. Every split therefore satisfies
\begin{equation}\label{eq:charge}
 w(I)<(c/2)^p b(I),\qquad d=2p-1.
\end{equation}
Let $T$ consist of internal nodes with no internal children. These intervals
are disjoint. Assign each internal node to a terminal internal descendant.
Each child has size at most $2/3$ of its parent, also when sizes are odd.
Nodes assigned to a fixed terminal node lie on its ancestor path at distinct
depths, so their weights sum to at most $w(J)/(1-r)$. Consequently
\[
 \sum_{I\text{ internal}}w(I)
 \le\frac{(c/2)^p}{1-r}\sum_{x_1}\alpha(x_1)^{-p}.
\]
Each child also has size at least $m/3$. Telescoping leaf weights over all
splits gives
\[
 \sum_{I\text{ leaf}}w(I)
 \le n^{-d}+2\,3^d\sum_{I\text{ internal}}w(I).
\]
The exact transverse cost of a leaf is $\lceil n/m\rceil^d\le(2n/m)^d$.
Multiply the last display by $(2n)^d$ and use
$\sum_Xa^{-p}=n^d\sum_{x_1}\alpha^{-p}$ to obtain~\eqref{eq:count}.
Finally $\sum_Xa^{-p}\le MN/\mu^p$ and
$\mu^{-p}\le2^p(1+\mu)^{-p}$ prove the stated target bound.
This completes the proof of Theorem~\ref{thm:laminate} without a logarithmic
charge at each tree level.

\subsection{A thin weak layer beyond minimum certificates}
Take integer $t\ge2$, $n=t^D$, $\mu=n^2$, $b=n^{2-2/D}$, and
\[
 \alpha(0)=b,\qquad \alpha(i)=H=\frac{n\mu-b}{n-1}\quad(i>0).
\]
Then $1\le\alpha\le2\mu$, its exact mean is $\mu$, and
\begin{equation}\label{eq:moment-layer}
 \operatorname{mean}((\mu/a)^p)
 =1+(1-1/n)(\mu/H)^p<2.
\end{equation}
The uniform bound for this unbounded sequence is $M=2$; the moments tend
to two, so a single fixed $M<2$ is not asserted. Moreover
\[
 nR_{\rm root}=n^{-1+2/D}
       +\frac{(n-1)^2}{n^2-b/n}<2.
\]
Thus at $c=1$ the algorithm accepts the root and uses one mean.
Yet $b\to\infty$ and $b/\mu\to0$, excluding every fixed coefficient-gap
condition $a\le B$ or $a\ge\theta\mu$ for fixed $B,\theta>0$ eventually.

Consider instead any partition whose cells pass the minimum-only sufficient
test $(\min_Ra)\lambda_1^0(R)\ge1$, where $\lambda_1^0$ is the actual
unweighted internal gap. For a singleton the mean-zero condition is vacuous.
Every cell meeting the weak layer has minimum $b$. A linear path test gives
$\lambda_1^0\le12/\ell^2$ in a cut coordinate of length $\ell\ge2$;
a Fourier test gives $\lambda_1^0<40/\ell^2$ in a full cycle coordinate.
Thus each transverse side of such a cell is at most $\sqrt{40b}$.
Covering the $n^d$ weak-layer vertices forces
\[
 \#\Pcal\ge40^{-d/2}n^{d/D}.
\]
This is a limitation of that certificate, since the preceding one-cell
construction works on the same coefficients.

\section{The cell cost of strong Green moments}
Now take homogeneous $a=\mu=n^2$. These original coefficient data have
$M=1$ and target order one. On a rectangle $R$ use cut coordinate paths,
or retain a periodic edge only when its coordinate spans the full period.
Let $L_R$ be its actual weighted Laplacian and $G_R=L_R^\dagger$, zero
on the constant kernel. For $s>0$ let
\[
 T_s(R)=\tr(G_R^s)=\sum_{\lambda_j(L_R)>0}\lambda_j(L_R)^{-s}.
\]
For a singleton set $T_s=0$. The strong certificate $T_s(R)\le1$ indeed
implies $\Var_R(u)\le\E_{a,R}(u)$, because its first positive inverse
eigenvalue raised to $s$ is one summand. Define $q_s(n)$ as the minimum
number of cells in any disjoint rectangular partition passing this test.
No hierarchy, aspect ratio or placement restriction is imposed.

\begin{theorem}\label{thm:green}
For fixed $D\ge3$ and $s>0$,
\[
 q_s(n)=
 \begin{cases}
 \Theta_{D,s}(n^{D-2s}),&0<s<D/2,\\
 \Theta_D(\log n),&s=D/2,\\
 \Theta_{D,s}(1),&s>D/2.
 \end{cases}
\]
The actual gap and a weak critical counting certificate both allow one
root cell on this family. This phase diagram concerns sufficient
certificate cost, not necessary observation rank.
\end{theorem}

\subsection{Finite spectra and all-partition lower bounds}
The cut path of length $\ell$ has eigenvectors
$\cos(\pi k(x+1/2)/\ell)$ and eigenvalues
$4\sin^2(\pi k/(2\ell))$, $0\le k<\ell$. The endpoint identity follows
from cosine addition and $\sin(\pi k)=0$. Distinct eigenvalues give a full
orthogonal basis. A full cycle has the usual Fourier eigenvalues
$4\sin^2(\pi k/n)$. At $n=2$ retain both positive-start periodic edges;
its nonzero eigenvalue is four, while the cut path eigenvalue is two.
Tensor products give the cut-rectangle spectrum
\begin{equation}\label{eq:spectrum}
 \lambda_k=4n^2\sum_{i=1}^D\sin^2\frac{\pi k_i}{2\ell_i},
 \qquad 0\le k_i<\ell_i,\quad k\ne0.
\end{equation}
Adding full-period edges increases ordered eigenvalues by min-max and keeps
the constant kernel. Comparing eigenvalues suffices for every $s>0$;
we do not assume an operator-power order valid for arbitrary powers.

Every such internal graph has $\lambda_{\max}\le4Dn^2$. Hence, with
$V=|R|$, $T_s(R)\ge(V-1)/(4Dn^2)^s$. A passing partition satisfies
\begin{equation}\label{eq:subcrit}
 q\ge\frac{n^D}{1+(4Dn^2)^s}.
\end{equation}
This proves the subcritical lower bound for arbitrary shapes.

For the endpoint set $p=D/2$, $m=\min_i\ell_i$,
$J(m)=\max(0,\lfloor\log_2m\rfloor-2)$ and
$A_D=2^{-D}(256D)^{-\lceil D/2\rceil}$. We claim
\begin{equation}\label{eq:blocks}
 T_p(R)\ge A_D(V/n^D)J(m).
\end{equation}
It is vacuous if $m<8$. Otherwise for
$j=0,\ldots,\lfloor\log_2m\rfloor-3$ put $t_j=2^{j+1}/m$ and select
the disjoint mode blocks
$\lceil t_j\ell_i\rceil\le k_i<\lceil2t_j\ell_i\rceil$ for every $i$.
Their sizes are at least $Vt_j^D/2^D$: for $x=t_j\ell_i\ge2$,
$\lceil2x\rceil-\lceil x\rceil\ge x/2$.
All selected frequencies have $k_i/\ell_i<2t_j\le1/2$ and satisfy
$\lambda_k<256Dn^2t_j^2$, under both path and retained-cycle spectra.
Their contribution is at least $A_DV/n^D$ per block. This proves
\eqref{eq:blocks}, including arbitrary aspect ratios.

In a partition with $q$ cells, those of volume at least $N/(2q)$ cover
at least half of $X$. Each of their sides is at least $n/(2q)$.
Summing~\eqref{eq:blocks} over these cells and using $T_p\le1$ gives
\[
 q\ge\frac{A_D}{2}
 \max\{0,\lfloor\log_2(n/(2q))\rfloor-2\}.
\]
For $n\ge2^{16}$ this implies $q\ge(A_D/8)\log_2 n$.
Indeed if $q\ge\sqrt n$ this is immediate; otherwise the expression in
braces is at least $(\log_2 n)/2-4\ge(\log_2 n)/4$.

\subsection{Matching actual meshes}
Use a Cartesian mesh with all side lengths at most an integer $M\le n$,
allowing shorter last intervals. From~\eqref{eq:spectrum} and
$\sin t\ge2t/\pi$ on $[0,\pi/2]$, group modes by $j=\max_i k_i$ to get
\begin{equation}\label{eq:upper}
 T_s(R)\le D2^{D-1}4^{-s}(M/n)^{2s}
          \sum_{j=1}^{M-1}j^{D-1-2s}.
\end{equation}
There are at most $D(j+1)^{D-1}\le D2^{D-1}j^{D-1}$ modes in the group.
When $2s<D$, the sum is at most a constant depending on $D,s$ times
$M^{D-2s}$. Choose $M=\max(1,\lfloor\varepsilon n^{2s/D}\rfloor)$
with fixed sufficiently small $\varepsilon>0$. Then every trace passes,
and the cell count is $O_{D,s}(n^{D-2s})$. Floor-zero cases use singletons.
When $2s=D$, the sum is at most $1+\log M$. Choosing
$M=\max(1,\lfloor\varepsilon n/(1+\log n)^{1/D}\rfloor)$ gives
$O_D(1+\log n)$ cells. When $2s>D$, the sum is at most
$1+(2s-D)^{-1}$; choose $M=\max(1,\lfloor\varepsilon n\rfloor)$
to get a constant cell count. Retained full-period edges only decrease
the traces. These constructions and the lower bounds prove
Theorem~\ref{thm:green}.

\subsection{Successful weak endpoint and exact finite distinctions}
For the cut root, $\lambda_k\ge4|k|^2$, so its true positive gap is at
least four. Define
\[
 W_p=\sup_{\lambda>0}\lambda^{-p}
          \#\{0<\lambda_j(L_X)<\lambda\}.
\]
There are no such eigenvalues for $\lambda<4$; for $\lambda\ge4$ their
number is at most $(1+\sqrt\lambda/2)^D\le\lambda^p$. Thus $W_p\le1$.
The same holds after adding periodic edges. Integrality of this counting
function excludes any eigenvalue below one, so this weak criterion really
certifies the successful one-mean inequality. The weak/strong ideal
distinction is classical~\cite{frank}; its use does not solve the original
coefficient-to-rectangle construction.

For $D=4,n=4$, the periodic coordinate spectrum at $a=16$ is
$0,32,64,32$. The root eigenvalues $32m$ have multiplicities
$\binom8m$. Its gap is32, but
$T_1(X)\ge56/96+70/128=217/192>1$; sixteen cut length-two boxes instead
have $T_1=103/384$ each. At the critical endpoint an exact separate
example takes $n=2$, $\mu=4$, $c=16$ and test $16^pT_p(R)\le1$.
A rectangle with $r$ full coordinates has spectrum $16m$ with multiplicity
$\binom rm$, so this test passes exactly when $r\le1$. Therefore its
minimum cell count is $2^{D-1}$, whereas the actual one-mean inequality
works with $K=c=16$. For $D=4$ the rejected root value is $865/144$.
Changing a fixed $c$ changes the constants in the phase diagram, not its
orders. The ordinary identity
$\sum_{x<y}R_{\rm eff}(x,y)=|R|\tr(G_R)$ follows by expanding
$(e_x-e_y)^TG_R(e_x-e_y)$ and using $G_R\mathbf1=0$; it only identifies
the all-pairs resistance sum with the $s=1$ test, not every resistance method.

\section{A coexact falsifier for representative fibres}
Lamination cannot be replaced by sampling one favourable fibre. In
$D=3,n=4$ let $c=15/2$ and set $a=1$ on $\{0,1,2\}^3$, $a=64$
elsewhere. Its exact mean is $\mu=2395/64$, $1\le a\le2\mu$, and
\[
 \operatorname{mean}((\mu/a)^{3/2})
 <\frac{27}{64}\frac{1178}{5}+\frac{37}{64}
 =\frac{31991}{320}<100.
\]
Here $\mu<38$ and $\sqrt\mu<31/5$. The fibre through $(3,3)$ in the
transverse coordinates has all coefficients64 and would pass
$nR=1/4<2/c=4/15$, leading to a false one-cell conclusion.

Let $G=\delta_{(1,1,1)}$ and use backward differences to define
$U_1=D_2^-G$, $U_2=-D_1^-G$, $U_3=0$. This is the codifferential of
an actual antisymmetric plaquette potential, up to its global sign
convention; backward divergence vanishes by commuting differences.
Both nonzero components have zero mean and squared norm two, so
$\|U\|^2=4$. An adjacent opposite delta pair has forward-gradient
energy six in its own direction and four in each other direction.
The total for $U$ is28. Every nonzero forward-gradient starting vertex
lies in the weak cube, so the actual weighted energy is also28.
Consequently $\E_a(U)=28<30=c\|U\|^2$, while the global mean penalty
vanishes. This refutes the representative-fibre mechanism even on the
actual coexact domain. A different partition remains allowed by the
original problem, so that problem has not been refuted.

\section{Literature comparison and verification boundaries}
Weighted path Cauchy estimates and weighted Hardy theory are classical;
Muckenhoupt~\cite{muckenhoupt} is relevant context. Adaptive tree
approximation has an extensive literature, including
Binev--DeVore~\cite{binev}. Our tree charge is proved explicitly rather
than imported from a near-best approximation theorem. Publisher/author
metadata for those two references were checked; unavailable full-text
retrievals are not represented as complete comparisons.

Nochetto--Ot\'arola--Salgado~\cite{nos}, Lemma3.1 and Theorem5.7,
give constructive weighted Poincar\'e and anisotropic rectangular
interpolation estimates in Muckenhoupt-weighted Sobolev spaces.
Their constants depend on the weight characteristic and mesh geometry.
Their inspected results are not the finite counting-norm observation
theorem proved here, nor a proof of its uniform constants from only a
global inverse moment. Conversely this distinction is not proof of
originality: broader weighted approximation and adaptive-partition
literature still requires comparison. In particular no claim is made to
invent weighted Poincar\'e inequalities, spectral counting, adaptive trees,
or weak Schatten ideals. The weak/strong distinction is separately
attributed to the established spectral literature~\cite{frank}.

The universal mathematical arguments above are analytical, not Lean
verified. Source units184 and185 have exact rational finite checks at
pinned Git commits. Unit184 has four matching normal/optimized payloads
from two fresh Colab clones, six rejecting mutations, five critical matrix
controls and six mode-block cases. Unit185 has four matching payloads
from two fresh clones, six global matrix controls,72 balanced-tree cases,
five weak-layer instances and six rejected intentional errors. It checks
the actual coexact energy28 versus30. These are selected finite tests,
not an all-size computation or independent replication. Full stdout,
stderr, actual exit codes, hashes, failures and resource records are in
the source evidence directories. No scientific computation ran on Windows.

\subsection{Model assessments and review context}
This exact supplement PDF has no inherited numerical grade or stars.
It is a draft pending its own declared exact-PDF assessment, not a public
AIRR release, human peer-reviewed paper or editorial acceptance. Public
version8 of the existing dossier has a separate exact-PDF AI assessment;
its score does not transfer to these bytes.

The proof of unit185 was read once in full in an adversarial source audit;
the auditor found no blocking defect in the laminate construction or
critical charge and recorded the uniform $M=2$ distinction and dimension,
singleton and target-order qualifications incorporated here. This audit
has no numeric score. Interface: OpenAI Codex, configured gpt-6-astra;
the actual served backend and reasoning setting were unexposed. Shared
account, provider and historical chat context mean the audit is neither
blind nor organizationally independent. The full report is retained,
including objections; it is not replaced by this summary. It did not
review this newly typeset exact PDF. No assessment was rerun to improve
a grade, and no report is represented as human peer review.

\subsection{Reproducibility and remaining question}
The source proof for185 has SHA256
\begin{center}\small\texttt{\detokenize{cee3556c1ff9dd99a69b3d62b20f65288189}}\\
\texttt{\detokenize{26cca2734b2728b51ce9e4fdd407}}\end{center}
The corresponding executed commit is
\texttt{dfc2205dcfb8642bf5639c8433e4bdbff38075a3};
source184 was executed at
\texttt{bedf55c593994aeba2cf47454c2f608870ecb871}.
This supplement's source and rendered-file hashes belong in its separate
build manifest; a PDF cannot contain its own final hash. The research
evidence is preserved at the pinned programme repository~\cite{evidence}.

The sole unrestricted critical rectangular question remains open.
A proposed extension must derive genuine transverse or all-fibre control
from its coefficient hypotheses; one favourable fibre is insufficient.
The strong Green test alone is abandoned at and below the critical
exponent as a uniform target-order mechanism by Theorem~\ref{thm:green}.
These findings do not reopen previously stopped routes, change PR29's
20/41 and TermSource0 baseline, or identify a finite sequence bound with
a continuum or interacting Yang--Mills mass gap.

\begin{thebibliography}{9}\small\raggedright
\bibitem{programme} L. Eriksson, \emph{Coarse information and finite rank
obstructions for flat lattice Hodge forms}, existing AIRR dossier
ARR-2026-7EPM3V3DP796BBRR, public version8.
\url{https://airr.science/papers/ARR-2026-7EPM3V3DP796BBRR/versions/v8/}.
Working paper with AI assessment, not editorial acceptance.
\bibitem{muckenhoupt} B. Muckenhoupt, \emph{Hardy's inequality with weights},
Studia Mathematica44 (1972),31--38.
\url{https://doi.org/10.4064/sm-44-1-31-38}.
\bibitem{binev} P. Binev and R. DeVore, \emph{Fast computation in adaptive
tree approximation}, Numerische Mathematik97 (2004),193--217.
\url{https://doi.org/10.1007/s00211-003-0493-6}.
\bibitem{nos} R. H. Nochetto, E. Ot\'arola and A. J. Salgado,
\emph{Piecewise polynomial interpolation in Muckenhoupt weighted Sobolev
spaces and applications}, arXiv:1402.1916v2 (2014), Lemma3.1,
Section5.2 and Theorem5.7. \url{https://arxiv.org/abs/1402.1916v2}.
\bibitem{frank} R. L. Frank, \emph{Cwikel's theorem and the CLR inequality},
arXiv:1206.3325v1 (2012), Section2.
\url{https://arxiv.org/abs/1206.3325v1}.
\bibitem{evidence} L. Eriksson programme research evidence, units184--185,
pinned at commit506004b3e7251aa881b08dde1a1bb14aff33e76f.
\href{https://github.com/lluiseriksson/THE-ERIKSSON-PROGRAMME/tree/506004b3e7251aa881b08dde1a1bb14aff33e76f/papers/flat-coarse-rank/research}{Source proofs and complete finite evidence}.
\end{thebibliography}
\end{document}
